% =========================================================
% Proof: Preservation of Convergence and Cauchy Sequences
% Source: notes/isometries/notes-isometries.tex
% Difficulty: Easy
% =========================================================

\subsection*{Preservation of Convergence and Cauchy Sequences}
\label{prf:isometry-preserves-convergence}

\begin{remark*}[Return]
$\leftarrow$ \hyperref[prop:isometry-preserves-convergence]{Back to Proposition in Notes}
\end{remark*}

\begin{proof}
\Claim If $f : X \to Y$ is a bijective isometry then:
(i) $x_n \to x$ in $X$ iff $f(x_n) \to f(x)$ in $Y$;
(ii) $(x_n)$ is Cauchy in $X$ iff $(f(x_n))$ is Cauchy in $Y$.

\Given A bijective isometry $f$ and a sequence $(x_n)$ in $X$.

\Goal To prove both biconditionals.

\Strategy Both statements reduce to the same observation:
$d_Y(f(x_n),f(x)) = d_X(x_n,x)$ for all $n$,
so the sequence of distances is identical in both spaces.

% -------------------------------------------------------
% YOUR PROOF HERE
% -------------------------------------------------------

\end{proof}

\begin{remark*}[Proof shape]
One-line observation: the isometry condition makes the distance sequences
in $X$ and $Y$ literally equal. Both convergence and Cauchy are purely
statements about those distances.
\end{remark*}

\begin{remark*}[Dependencies]
\begin{itemize}
  \item Definition~\ref{def:isometry}: $d_Y(f(x),f(y)) = d_X(x,y)$.
  \item Lemma~\ref{lem:convergence-preserved}: convergence preserved (one direction,
        proved earlier).
  \item Definition~\ref{def:cauchy}: Cauchy sequence.
\end{itemize}
\end{remark*}
